\section{Projective limits in homological algebra}
\label{section:0.13}

\subsection{The Mittag-Leffler condition}
\label{subsection:0.13.1}

\begin{env}[13.1.1]
\label{0.13.1.1}
Let $\cat{C}$ be an abelian category in which infinite products exist
(axiom AB~3* of T, 1.5).
Then the infimum of a family of subobjects of an object of $\cat{C}$
exists, and every projective system of objects of $\cat{C}$ has a
projective limit; the latter is a left exact functor of the projective
system in question (T, 1.8).
\oldpage[0\textsubscript{III}]{65}
Let $(A_\alpha,f_{\alpha\beta})$ be a projective system of objects of
$\cat{C}$ whose index set $\mathrm{I}$ is filtered to the right.
Put
\[
  A=\varprojlim A_\alpha,
\]
and, for each $\alpha\in\mathrm{I}$, let
$f_\alpha:A\longrightarrow A_\alpha$ be the canonical morphism.
For every $\alpha\in\mathrm{I}$, the subobjects
$f_{\alpha\beta}(A_\beta)$, for $\alpha\leq\beta$, form a decreasing
filtered family of subobjects of $A_\alpha$.
The subobject
\[
  A'_\alpha
  =
  \inf_{\beta\geq\alpha}f_{\alpha\beta}(A_\beta)
\]
is called the subobject of \emph{universal images} in $A_\alpha$.
It is clear that
\[
  f_\alpha(A)\subset A'_\alpha,
  \qquad
  f_{\alpha\beta}(A'_\beta)\subset A'_\alpha
  \quad(\alpha\leq\beta).
\]
Thus
$(A'_\alpha,f_{\alpha\beta}|A'_\beta)$ is a projective system and
\[
  A=\varprojlim A'_\alpha.
\]
\end{env}

\begin{env}[13.1.2]
\label{0.13.1.2}
For a projective system $(A_\alpha,f_{\alpha\beta})$ in $\cat{C}$,
the following is called the \emph{Mittag-Leffler condition}:
\begin{itemize}
  \item[\textbf{(ML)}]
  \emph{For every index $\alpha$, there is a $\beta\geq\alpha$ such
  that, for every $\gamma\geq\beta$,}
  \[
    f_{\alpha\gamma}(A_\gamma)
    =
    f_{\alpha\beta}(A_\beta).
  \]
\end{itemize}

It is clear that condition (ML) holds if all the
$f_{\alpha\beta}$ are epimorphisms.
Conversely, suppose that (ML) holds, and for each
$\alpha\in\mathrm{I}$ let $A'_\alpha$ be the subobject of universal
images in $A_\alpha$.
Then, for $\alpha\leq\beta$, the restriction of
$f_{\alpha\beta}$ to $A'_\beta$ is an \emph{epimorphism}
\[
  A'_\beta\longrightarrow A'_\alpha.
\]
Indeed, choose $\gamma\geq\beta$ such that
\[
  f_{\beta\delta}(A_\delta)
  =
  f_{\beta\gamma}(A_\gamma)
  \qquad(\delta\geq\gamma).
\]
Then $A'_\beta=f_{\beta\gamma}(A_\gamma)$, and therefore
\[
  f_{\alpha\beta}(A'_\beta)
  =
  f_{\alpha\gamma}(A_\gamma),
  \qquad
  A'_\alpha
  =
  f_{\alpha\gamma}(A_\gamma)
  =
  f_{\alpha\beta}(A'_\beta).
\]

Condition (ML) also holds when the objects $A_\alpha$ are
\emph{Artinian} in $\cat{C}$, that is, when every family of
subobjects of $A_\alpha$ has a minimal member.
Indeed, a minimal member of the decreasing filtered family
$\bigl(f_{\alpha\beta}(A_\beta)\bigr)$ of subobjects of $A_\alpha$
is necessarily the smallest of those subobjects.
\end{env}

\begin{env}[13.1.3]
\label{0.13.1.3}
\emph{Remark} (13.1.3).
Condition (ML) can also be formulated when $\cat{C}$ is, for
example, the category of sets.
One can again define the subset of universal images in $A_\alpha$,
and the observations made in \sref{0.13.1.1} and
\sref{0.13.1.2} remain valid.
\end{env}

\subsection{The Mittag-Leffler condition for abelian groups}
\label{subsection:0.13.2}

\begin{proposition}[13.2.1]
\label{0.13.2.1}
Let
\[
  0\longrightarrow A_\alpha
  \xrightarrow{u_\alpha}B_\alpha
  \xrightarrow{v_\alpha}C_\alpha
  \longrightarrow0
\]
be an exact sequence of projective systems of abelian groups, relative
to the same filtered index set $\mathrm{I}$.
\begin{enumerate}
  \item[{\rm(i)}]
  If $(B_\alpha)$ satisfies {\rm(ML)}, then so does $(C_\alpha)$.
  \item[{\rm(ii)}]
  If $(A_\alpha)$ and $(C_\alpha)$ satisfy {\rm(ML)}, then so does
  $(B_\alpha)$.
\end{enumerate}
\end{proposition}

\begin{proof}
Let $(f_{\alpha\beta})$, $(g_{\alpha\beta})$, and
$(h_{\alpha\beta})$ be the systems of homomorphisms defining
$(A_\alpha)$, $(B_\alpha)$, and $(C_\alpha)$, respectively.

(i) Suppose that
\[
  g_{\alpha\beta}(B_\beta)
  =
  g_{\alpha\lambda}(B_\lambda)
  \qquad(\lambda\geq\beta).
\]
Since $v_\beta$ and $v_\lambda$ are surjective,
\[
  h_{\alpha\beta}(C_\beta)
  =
  v_\alpha\bigl(g_{\alpha\beta}(B_\beta)\bigr)
  =
  v_\alpha\bigl(g_{\alpha\lambda}(B_\lambda)\bigr)
  =
  h_{\alpha\lambda}(C_\lambda)
\]
for $\lambda\geq\beta$.

(ii) Let $\alpha\in\mathrm{I}$.
Choose $\beta\geq\alpha$ such that
\[
  f_{\alpha\beta}(A_\beta)
  =
  f_{\alpha\lambda}(A_\lambda)
  \qquad(\lambda\geq\beta),
\]
and then choose $\gamma\geq\beta$ such that
\[
  h_{\beta\gamma}(C_\gamma)
  =
  h_{\beta\lambda}(C_\lambda)
  \qquad(\lambda\geq\gamma).
\]
Let $y_\alpha\in g_{\alpha\gamma}(B_\gamma)$, write
$y_\alpha=g_{\alpha\gamma}(y_\gamma)$ with $y_\gamma\in B_\gamma$,
and put $y_\beta=g_{\beta\gamma}(y_\gamma)$.
Then
\[
  v_\beta(y_\beta)
  =
  h_{\beta\gamma}\bigl(v_\gamma(y_\gamma)\bigr).
\]
For every $\lambda\geq\gamma$, there is, by hypothesis, a
$y_\lambda\in B_\lambda$ such that
\[
  h_{\beta\gamma}\bigl(v_\gamma(y_\gamma)\bigr)
  =
  h_{\beta\lambda}\bigl(v_\lambda(y_\lambda)\bigr).
\]
Hence
$v_\beta(y_\beta-g_{\beta\lambda}(y_\lambda))=0$, and therefore
there is an $x_\beta\in A_\beta$
\oldpage[0\textsubscript{III}]{66}
such that
\[
  y_\beta
  =
  g_{\beta\lambda}(y_\lambda)+u_\beta(x_\beta).
\]
It follows that
\[
  y_\alpha
  =
  g_{\alpha\lambda}(y_\lambda)
  +
  u_\alpha\bigl(f_{\alpha\beta}(x_\beta)\bigr).
\]
Since $\lambda\geq\beta$, there is an $x_\lambda\in A_\lambda$ such
that
\[
  f_{\alpha\beta}(x_\beta)=f_{\alpha\lambda}(x_\lambda).
\]
Consequently,
\[
  y_\alpha
  =
  g_{\alpha\lambda}
  \bigl(y_\lambda+u_\lambda(x_\lambda)\bigr)
  \in g_{\alpha\lambda}(B_\lambda),
\]
which proves the assertion.
\end{proof}

\begin{proposition}[13.2.2]
\label{0.13.2.2}
Let $\mathrm{I}$ be a filtered ordered set having a countable
cofinal subset, and let
\[
  0\longrightarrow A_\alpha
  \xrightarrow{u_\alpha}B_\alpha
  \xrightarrow{v_\alpha}C_\alpha
  \longrightarrow0
\]
be an exact sequence of projective systems of abelian groups indexed
by $\mathrm{I}$.
If $(A_\alpha)$ satisfies {\rm(ML)}, then the sequence
\[
  0\longrightarrow\varprojlim A_\alpha
  \longrightarrow\varprojlim B_\alpha
  \longrightarrow\varprojlim C_\alpha
  \longrightarrow0
\]
is exact.
\end{proposition}

\begin{proof}
It remains only to prove that the homomorphism
\[
  v=\varprojlim v_\alpha:
  \varprojlim B_\alpha\longrightarrow\varprojlim C_\alpha
\]
is surjective.
Let $z=(z_\alpha)\in\varprojlim C_\alpha$, and put
\[
  E_\alpha=v_\alpha^{-1}(z_\alpha).
\]
The $E_\alpha$ plainly form a projective system of nonempty sets under
the restrictions of the homomorphisms
$g_{\alpha\beta}:B_\beta\longrightarrow B_\alpha$.

We show that this projective system satisfies (ML).
Identifying $A_\alpha$ with a subgroup of $B_\alpha$ by means of
$u_\alpha$, choose, for every $\alpha\in\mathrm{I}$, an index
$\beta\geq\alpha$ such that
\[
  g_{\alpha\beta}(A_\beta)
  =
  g_{\alpha\lambda}(A_\lambda)
  \qquad(\lambda\geq\beta).
\]
We claim that also
\[
  g_{\alpha\beta}(E_\beta)
  =
  g_{\alpha\lambda}(E_\lambda)
  \qquad(\lambda\geq\beta).
\]
Indeed, take $y_\lambda\in E_\lambda$ and set
\[
  y_\beta=g_{\beta\lambda}(y_\lambda),
  \qquad
  y_\alpha=g_{\alpha\lambda}(y_\lambda).
\]
Let $y'_\alpha=g_{\alpha\beta}(y'_\beta)$ for some
$y'_\beta\in E_\beta$.
Then $y'_\beta-y_\beta=x_\beta\in A_\beta$, and by hypothesis there
is an $x_\lambda\in A_\lambda$ such that
\[
  g_{\alpha\beta}(x_\beta)
  =
  g_{\alpha\lambda}(x_\lambda).
\]
Thus
\[
  y'_\alpha
  =
  g_{\alpha\beta}(y_\beta)+g_{\alpha\beta}(x_\beta)
  =
  g_{\alpha\lambda}(y_\lambda+x_\lambda)
  \in g_{\alpha\lambda}(E_\lambda),
\]
which proves the claim.

It is known (Bourbaki, \emph{Top.\ g\'{e}n.}, chap.~II,
\textsection3, th.~1) that, under the hypotheses imposed on
$\mathrm{I}$, a projective system of nonempty sets satisfying (ML)
has a nonempty projective limit.
There is therefore a point
$y=(y_\alpha)\in\varprojlim E_\alpha$.
Since $v_\alpha(y_\alpha)=z_\alpha$ for every $\alpha$, one has
$z=v(y)$.
\end{proof}

\begin{proposition}[13.2.3]
\label{0.13.2.3}
Under the hypotheses on $\mathrm{I}$ of \sref{0.13.2.2}, let
$(\mathrm{K}^\bullet_\alpha)_{\alpha\in\mathrm{I}}$ be a projective
system of complexes of abelian groups
\[
  \mathrm{K}^\bullet_\alpha
  =
  (\mathrm{K}^n_\alpha)_{n\in\bb{Z}},
\]
whose differentials have degree $+1$.
For every $n$, there is a canonical functorial homomorphism
\begin{equation}
\label{0.13.2.3.1}
\tag{13.2.3.1}
  h_n:
  \HH^n\bigl(\varprojlim\mathrm{K}^\bullet_\alpha\bigr)
  \longrightarrow
  \varprojlim\HH^n(\mathrm{K}^\bullet_\alpha).
\end{equation}
If, for every degree $n$, the projective system of abelian groups
$(\mathrm{K}^n_\alpha)_{\alpha\in\mathrm{I}}$ satisfies {\rm(ML)},
then every $h_n$ is surjective.
If, in addition, the projective system
$(\HH^{n-1}(\mathrm{K}^\bullet_\alpha))_{\alpha\in\mathrm{I}}$
satisfies {\rm(ML)}, then $h_n$ is bijective.
\end{proposition}

\begin{proof}
Put
\[
  \mathrm{K}^n=\varprojlim\mathrm{K}^n_\alpha
  \qquad(n\in\bb{Z}).
\]
The homomorphisms $h_n$ are defined by the commutative diagrams
\[
  \xymatrix{
    \cdots\ar[r]
      & \mathrm{K}^{n-1}\ar[r]\ar[d]
      & \mathrm{K}^n\ar[r]\ar[d]
      & \mathrm{K}^{n+1}\ar[r]\ar[d]
      & \cdots\\
    \cdots\ar[r]
      & \mathrm{K}^{n-1}_\alpha\ar[r]
      & \mathrm{K}^n_\alpha\ar[r]
      & \mathrm{K}^{n+1}_\alpha\ar[r]
      & \cdots,
  }
\]
\oldpage[0\textsubscript{III}]{67}
the differentials in $\mathrm{K}^\bullet$ being the projective limits
of the corresponding differentials in the
$\mathrm{K}^\bullet_\alpha$.
Consider the exact sequences
\[
  (*_n)\qquad
  0\longrightarrow
  \mathrm{B}^n(\mathrm{K}^\bullet_\alpha)
  \longrightarrow
  \mathrm{Z}^n(\mathrm{K}^\bullet_\alpha)
  \longrightarrow
  \HH^n(\mathrm{K}^\bullet_\alpha)
  \longrightarrow0
\]
and
\[
  (**_n)\qquad
  0\longrightarrow
  \mathrm{Z}^{n-1}(\mathrm{K}^\bullet_\alpha)
  \longrightarrow
  \mathrm{K}^{n-1}_\alpha
  \longrightarrow
  \mathrm{B}^n(\mathrm{K}^\bullet_\alpha)
  \longrightarrow0.
\]
The hypothesis and Proposition~\sref{0.13.2.1}~(i) show that the
projective system
$(\mathrm{B}^n(\mathrm{K}^\bullet_\alpha))_{\alpha\in\mathrm{I}}$
satisfies (ML) for every $n$.
It follows from \sref{0.13.2.2} that
\[
  (***_n)\qquad
  0\longrightarrow
  \varprojlim_\alpha\mathrm{B}^n(\mathrm{K}^\bullet_\alpha)
  \longrightarrow
  \varprojlim_\alpha\mathrm{Z}^n(\mathrm{K}^\bullet_\alpha)
  \longrightarrow
  \varprojlim_\alpha\HH^n(\mathrm{K}^\bullet_\alpha)
  \longrightarrow0
\]
is exact.
Now
$\varprojlim_\alpha\mathrm{B}^n(\mathrm{K}^\bullet_\alpha)$
identifies with a subgroup of $\mathrm{K}^{n+1}$ containing
$\mathrm{B}^n(\mathrm{K}^\bullet)$, while
$\varprojlim_\alpha\mathrm{Z}^n(\mathrm{K}^\bullet_\alpha)$
identifies with a subgroup of $\mathrm{Z}^n(\mathrm{K}^\bullet)$.
Consequently, $h_n$ is surjective.

Suppose now, in addition, that
$(\HH^{n-1}(\mathrm{K}^\bullet_\alpha))_{\alpha\in\mathrm{I}}$
satisfies (ML).
The exact sequences $(*_{n-1})$ and
Proposition~\sref{0.13.2.1}~(ii) show that
$(\mathrm{Z}^{n-1}(\mathrm{K}^\bullet_\alpha))_{\alpha\in\mathrm{I}}$
satisfies (ML).
Applying \sref{0.13.2.2} to the exact sequences $(**_n)$ gives an
exact sequence
\[
  0\longrightarrow
  \varprojlim_\alpha
    \mathrm{Z}^{n-1}(\mathrm{K}^\bullet_\alpha)
  \longrightarrow
  \mathrm{K}^{n-1}
  \xrightarrow{u}
  \varprojlim_\alpha
    \mathrm{B}^n(\mathrm{K}^\bullet_\alpha)
  \longrightarrow0.
\]
Since
$\varprojlim_\alpha\mathrm{B}^n(\mathrm{K}^\bullet_\alpha)
\supset\mathrm{B}^n(\mathrm{K}^\bullet)$
and the composite of $u$ with the injection
\[
  \varprojlim_\alpha\mathrm{B}^n(\mathrm{K}^\bullet_\alpha)
  \longrightarrow\mathrm{K}^n
\]
is the differential
$\mathrm{K}^{n-1}\longrightarrow\mathrm{K}^n$, the surjectivity of
$u$ implies
\[
  \varprojlim_\alpha\mathrm{B}^n(\mathrm{K}^\bullet_\alpha)
  =
  \mathrm{B}^n(\mathrm{K}^\bullet).
\]
Thus $h_n$ is injective.
\end{proof}

\begin{env}[13.2.4]
\label{0.13.2.4}
\emph{Remarks} (13.2.4).
\begin{enumerate}
  \item[{\rm(i)}]
  The argument of \sref{0.13.2.2}
  (cf.\ Bourbaki, \emph{loc.\ cit.}) shows that the conclusion of
  that proposition remains valid if one assumes only that the
  $A_\alpha$ can be equipped with structures of complete metrizable
  spaces in which translations are homeomorphisms, that the maps
  $f_{\alpha\beta}:A_\beta\longrightarrow A_\alpha$ defining the
  projective system are uniformly continuous for the metrics in
  question, and that $(A_\alpha)$ satisfies the condition
  \begin{itemize}
    \item[\textbf{(ML')}]
    \emph{For every index $\alpha$, there is a $\beta\geq\alpha$ such
    that, for every $\gamma\geq\beta$,
    $f_{\alpha\gamma}(A_\gamma)$ is dense in
    $f_{\alpha\beta}(A_\beta)$.}
  \end{itemize}

  This gives an analogous supplement to \sref{0.13.2.3}.
  Suppose that $\mathrm{K}^n_\alpha=0$ for $n<0$ and every $\alpha$;
  that
  $(\mathrm{K}^n_\alpha)_{\alpha\in\mathrm{I}}$ satisfies (ML) for
  $n\geq0$; and that
  $A_\alpha=\HH^0(\mathrm{K}^\bullet_\alpha)$ can be equipped with
  metric-space structures having the properties just stated.
  Then the conclusions of \sref{0.13.2.3} are unchanged for
  $n\geq2$, and $h_1$ is moreover \emph{bijective}.
  Indeed, the argument of \sref{0.13.2.2} also shows that
  $(\mathrm{B}^1(\mathrm{K}^\bullet_\alpha))_{\alpha\in\mathrm{I}}$
  satisfies (ML), that the sequence $(***_1)$ is exact, and finally
  that
  \[
    \varprojlim_\alpha
    \mathrm{B}^1(\mathrm{K}^\bullet_\alpha)
    =
    \mathrm{B}^1(\mathrm{K}^\bullet).
  \]
  This proves, among other things, the assertions of (T, 3.10.2).

  \item[{\rm(ii)}]
  One can introduce the right derived functors
  $\varprojlim^{(n)}$ of the functor $\varprojlim$, and obtain more
  complete statements than those above~\cite{III-28}.
\end{enumerate}
\end{env}

\subsection{Application: cohomology of a projective limit of sheaves}
\label{subsection:0.13.3}

\begin{proposition}[13.3.1]
\label{0.13.3.1}
\oldpage[0\textsubscript{III}]{68}
Let $X$ be a topological space, let
$(\sh{F}_k)_{k\in\bb{N}}$ be a projective system of sheaves of
abelian groups on $X$, and put
$\sh{F}=\varprojlim_k\sh{F}_k$.
Suppose that the following conditions hold:
\begin{enumerate}
  \item[{\rm(i)}]
  There is a basis $\mathfrak{B}$ for the topology of $X$ such that,
  for every $U\in\mathfrak{B}$ and every $i\geq0$, the projective
  system
  $(\HH^i(U,\sh{F}_k))_{k\in\bb{N}}$ satisfies (ML).

  \item[{\rm(ii)}]
  For every $x\in X$ and every $i>0$, one has
  \[
    \varinjlim_U\bigl(\varprojlim_k
    \HH^i(U,\sh{F}_k)\bigr)=0,
  \]
  where $U$ ranges over the neighbourhoods of $x$ belonging to
  $\mathfrak{B}$.

  \item[{\rm(iii)}]
  The homomorphisms
  $u_{hk}:\sh{F}_k\longrightarrow\sh{F}_h$ $(h\leq k)$ defining the
  projective system $(\sh{F}_k)$ are surjective.
\end{enumerate}
Under these conditions, for every $i>0$, the canonical homomorphism
\[
  h_i:\HH^i(X,\sh{F})
  \longrightarrow
  \varprojlim_k\HH^i(X,\sh{F}_k)
\]
is surjective. If, moreover, for a given value of $i$ the projective
system
$(\HH^{i-1}(X,\sh{F}_k))_{k\in\bb{N}}$ satisfies (ML), then $h_i$ is
bijective.
\end{proposition}

\begin{proof}
\emph{a)}
We first suppose that the $\sh{F}_k$, as well as the kernels
$\sh{N}_{hk}$ of the $u_{hk}$, are flasque; we shall then show that
condition~{\rm(iii)} in the statement implies
$\HH^i(X,\sh{F})=0$ for $i>0$.
It suffices to prove that, for every open subset $U$ of $X$ and every
covering $\mathfrak{U}$ of $U$ by open subsets of $U$, one has
\[
  \CHH^i(\mathfrak{U},\sh{F})=0
  \qquad (i>0).
\]
Indeed, it then follows first that, for \v{C}ech cohomology,
$\check{\HH}^{\,i}(U,\sh{F})=0$ for every $i>0$, and hence, by
(G, II, 5.9.2) applied to the set of all open subsets of $X$, that
$\HH^i(X,\sh{F})=0$ for every $i>0$.

Since the $\sh{F}_k$ are flasque, one has
$\CHH^i(\mathfrak{U},\sh{F}_k)=0$ for every $i>0$
(G, II, 5.2.3).
For each $k$, consider the complex
$\check{\mathrm{C}}^\bullet(\mathfrak{U},\sh{F}_k)$ of alternating
cochains on the nerve of the covering $\mathfrak{U}$ (G, II, 5.1);
these complexes plainly form a projective system of complexes of
abelian groups.
Let us show that all the maps
\[
  \check{\mathrm{C}}^\bullet(\mathfrak{U},\sh{F}_k)
  \longrightarrow
  \check{\mathrm{C}}^\bullet(\mathfrak{U},\sh{F}_h)
  \qquad(h\leq k)
\]
are surjective.
By definition, it is enough to show that, for every open subset $V$
of $X$, the map
$\Gamma(V,\sh{F}_k)\longrightarrow\Gamma(V,\sh{F}_h)$ is
surjective. But the sequence
\[
  0\longrightarrow\sh{N}_{hk}\longrightarrow\sh{F}_k
  \longrightarrow\sh{F}_h\longrightarrow0
\]
is exact by hypothesis and therefore gives the exact cohomology
sequence
\[
  \Gamma(V,\sh{F}_k)\longrightarrow\Gamma(V,\sh{F}_h)
  \longrightarrow\HH^1(V,\sh{N}_{hk})=0,
\]
since $\sh{N}_{hk}$ is flasque.

\oldpage[0\textsubscript{III}]{69}
The projective system
$(\check{\mathrm{C}}^\bullet(\mathfrak{U},\sh{F}_k))_{k\in\bb{N}}$
therefore satisfies (ML), and the same is true of
$(\CHH^i(\mathfrak{U},\sh{F}_k))_{k\in\bb{N}}$ for every $i\geq0$:
for $i>0$ this is trivial, while for $i=0$ it follows from what
precedes, since
$\CHH^0(\mathfrak{U},\sh{F}_k)=\Gamma(U,\sh{F}_k)$
(G, II, 5.2.2).
We may therefore apply \sref{0.13.2.3}, which gives
\[
  \CHH^i(\mathfrak{U},\sh{F})
  =
  \varprojlim_k\CHH^i(\mathfrak{U},\sh{F}_k)
  =0
  \qquad(i>0).
\]

\emph{b)}
We now pass to the general case. For each $k\in\bb{N}$, consider the
\emph{canonical resolution}
\[
  \sh{C}^\bullet(X,\sh{F}_k)
  =
  \bigl(\sh{C}^i(X,\sh{F}_k)\bigr)_{i\geq0}
\]
of $\sh{F}_k$ by flasque sheaves (G, II, 4.3).
For every $i\geq0$, it is clear that
$(\sh{C}^i(X,\sh{F}_k))_{k\in\bb{N}}$ is a projective system of
flasque sheaves; let us show that it satisfies the conditions of
part~\emph{a)}.
Indeed, if $\sh{N}_{hk}$ denotes the kernel of $u_{hk}$ for $h\leq k$,
then the sequence
\[
  0\longrightarrow\sh{N}_{hk}\longrightarrow\sh{F}_k
  \longrightarrow\sh{F}_h\longrightarrow0
\]
is exact by~{\rm(iii)}, and our assertion follows from the exactness
of the functor
$\sh{A}\longmapsto\sh{C}^i(X,\sh{A})$ (G, II, 4.3).

Put
\[
  \sh{G}^i=\varprojlim_k\sh{C}^i(X,\sh{F}_k).
\]
By part~\emph{a)}, one has
$\HH^j(X,\sh{G}^i)=0$ for $j>0$ and $i\geq0$.
We shall show that
$\sh{G}^\bullet=(\sh{G}^i)_{i\geq0}$ is a resolution of the sheaf
$\sh{F}$. Since $\HH^j(X,\sh{G}^i)=0$ for $j>0$, the cohomology
$\HH^\bullet(X,\sh{F})$ will then be equal to
$\HH^\bullet(\Gamma(X,\sh{G}^\bullet))$ (G, II, 4.7.1).

Taking the projective limit of the exact sequences
\[
  0\longrightarrow\sh{F}_k
  \longrightarrow\sh{C}^0(X,\sh{F}_k)
  \longrightarrow\sh{C}^1(X,\sh{F}_k)
  \longrightarrow\cdots
\]
plainly gives a complex of sheaves of abelian groups
\[
  0\longrightarrow\sh{F}
  \longrightarrow\sh{G}^0
  \longrightarrow\sh{G}^1
  \longrightarrow\cdots.
\]
To prove our assertion, it remains to establish that
$\sh{H}^i(\sh{G}^\bullet)=0$ for $i>0$.
This sheaf is generated by the presheaf
\[
  U\longmapsto\HH^i\bigl(\Gamma(U,\sh{G}^\bullet)\bigr)
\]
(G, II, 4.1).
Now the complex $\Gamma(U,\sh{G}^\bullet)$ is the projective limit of
the projective system of complexes of abelian groups
\[
  \bigl(\Gamma(U,\sh{C}^\bullet(X,\sh{F}_k))\bigr)_{k\in\bb{N}}
\]
(0\textsubscript{I}, \hyperref[0.3.2.6]{3.2.6}), while the canonical resolution
$\sh{C}^\bullet(U,\sh{F}_k|U)$ is induced on $U$ by
$\sh{C}^\bullet(X,\sh{F}_k)$.
By hypothesis~{\rm(i)}, for every $U\in\mathfrak{B}$ we may apply
\sref{0.13.2.3} to this projective system of complexes, and hence
\[
  \HH^i\bigl(\Gamma(U,\sh{G}^\bullet)\bigr)
  =
  \varprojlim_k\HH^i(U,\sh{F}_k)
  \qquad(i\geq0).
\]
Hypothesis~{\rm(ii)} then shows, by definition, that the sheaves
$\sh{H}^i(\sh{G}^\bullet)$ vanish for $i>0$.

Thus, for every $i\geq0$,
\[
  \HH^i(X,\sh{F})
  =
  \HH^i\bigl(\Gamma(X,\sh{G}^\bullet)\bigr),
  \qquad
  \Gamma(X,\sh{G}^\bullet)
  =
  \varprojlim_k\Gamma(X,\sh{C}^\bullet(X,\sh{F}_k)).
\]
We have just observed that all the maps
\[
  \Gamma(X,\sh{C}^i(X,\sh{F}_k))
  \longrightarrow
  \Gamma(X,\sh{C}^i(X,\sh{F}_h))
  \qquad(h\leq k)
\]
are surjective; the conclusion therefore follows once more from
\sref{0.13.2.3}.
\end{proof}

\begin{env}[13.3.2]
\label{0.13.3.2}
\emph{Remarks} (13.3.2).
\begin{enumerate}
  \item[{\rm(i)}]
  The statement of \sref{0.13.3.1} is of interest only for $i>0$,
  since, for $i=0$, $h_i$ is always an isomorphism without any
  hypothesis (0\textsubscript{I}, \hyperref[0.3.2.6]{3.2.6}).

  \item[{\rm(ii)}]
  Conditions~{\rm(i)} and~{\rm(ii)} of \sref{0.13.3.1} are satisfied,
  in particular, if
  $\HH^i(U,\sh{F}_k)=0$ for every $k$, every $i>0$, and every
  $U\in\mathfrak{B}$, and if, for $U\in\mathfrak{B}$, the maps
  \[
    \Gamma(U,\sh{F}_k)\longrightarrow\Gamma(U,\sh{F}_h)
  \]
  are surjective. This is the most frequent case in which
  \sref{0.13.3.1} is applied.
\end{enumerate}
\end{env}

\subsection{Mittag-Leffler condition and graded objects associated with projective systems}
\label{subsection:0.13.4}
\label{0.13.4}

\begin{env}[13.4.1]
\label{0.13.4.1}
Let
$\bb{A}=(A_k,u_{kh})_{k\in\bb{Z}}$ be a projective system in an
abelian category $\cat{C}$.
We say that it is \emph{bounded below} if there is a $k_0$ such that
$A_k=0$ for $k<k_0$.
On each $A_k$ we define a \emph{filtration}
$(\mathrm{F}^p(A_k))_{p\in\bb{Z}}$ by
\begin{equation}
\label{0.13.4.1.1}
\tag{13.4.1.1}
  \mathrm{F}^p(A_k)=
  \begin{cases}
    \Ker(A_k\longrightarrow A_{p-1}) & \text{for }p\leq k+1,\\
    0 & \text{for }p\geq k+1.
  \end{cases}
\end{equation}

\oldpage[0\textsubscript{III}]{70}
By hypothesis,
$\mathrm{F}^{k_0}(A_k)=A_k$ and
$\mathrm{F}^{k+1}(A_k)=0$; in other words, the filtration under
consideration is finite (\sref{0.11.1.3}).
The graded objects associated with this filtration are therefore
\[
  \mathrm{gr}^p(A_k)
  =
  \Ker(A_k\longrightarrow A_{p-1})
  \big/
  \Ker(A_k\longrightarrow A_p).
\]
Consequently, $\mathrm{gr}^p(A_k)$ is isomorphic to the image under
$A_k\longrightarrow A_p$ of
$\Ker(A_k\longrightarrow A_{p-1})$.
By the transitivity of the morphisms defining a projective system,
one therefore has
\begin{equation}
\label{0.13.4.1.2}
\tag{13.4.1.2}
  \mathrm{gr}^p(A_k)
  =
  \Ker(A_p\longrightarrow A_{p-1})
  \cap
  \Im(A_k\longrightarrow A_p).
\end{equation}
But \sref{0.13.4.1.1} gives
$\Ker(A_p\longrightarrow A_{p-1})=\mathrm{gr}^p(A_p)$, and hence
\begin{equation}
\label{0.13.4.1.3}
\tag{13.4.1.3}
  \mathrm{gr}^p(A_k)
  =
  \mathrm{gr}^p(A_p)
  \cap
  \Im(A_k\longrightarrow A_p).
\end{equation}

The preceding definitions also show that, for $k\leq h$,
\[
  u_{kh}\bigl(\mathrm{F}^p(A_h)\bigr)
  \subset
  \mathrm{F}^p(A_k),
\]
and consequently, for every $p\in\bb{Z}$, the
$\mathrm{gr}^p(u_{kh})$ define a projective system
$(\mathrm{gr}^p(A_k))_{k\in\bb{Z}}$.
\end{env}

\begin{env}[13.4.2]
\label{0.13.4.2}
We say that the projective system $\bb{A}$ is
\emph{essentially constant} if the morphisms
$A_{k+1}\longrightarrow A_k$ are isomorphisms for all sufficiently
large $k$.
We say that $\bb{A}$ is \emph{strict} if the morphisms
$A_i\longrightarrow A_j$ $(j\leq i)$ are epimorphisms.

When $\bb{A}$ is strict, it follows from \sref{0.13.4.1.3} that, for
$p\leq k\leq h$, the canonical morphism
\[
  \mathrm{gr}^p(A_h)\longrightarrow\mathrm{gr}^p(A_k)
\]
is an isomorphism; in other words, the projective system
$(\mathrm{gr}^p(A_k))_{k\in\bb{Z}}$ is essentially constant.
The sequence of objects $\mathrm{gr}^p(A_p)$, identified for every
$p$ with $\varprojlim_k\mathrm{gr}^p(A_k)$, is then denoted
$\mathrm{gr}^\bullet(\bb{A})$ and is called the
\emph{graded object associated with the strict projective system}
$\bb{A}=(A_k)$.

Suppose now that the bounded-below projective system $\bb{A}$
satisfies (ML).
By \sref{0.13.1.2}, the projective system
$\bb{A}'=(A'_k)$ of ``universal image'' objects is strict; it is also
bounded below.
The graded object $\mathrm{gr}^\bullet(\bb{A}')$ associated with
$\bb{A}'$ is then likewise called the graded object
\emph{associated with} $\bb{A}$ and is denoted
$\mathrm{gr}^\bullet(\bb{A})$.
\end{env}

\begin{proposition}[13.4.3]
\label{0.13.4.3}
Let $\bb{A}=(A_k)_{k\in\bb{Z}}$ be a bounded-below projective system
in an abelian category $\cat{C}$.
The following conditions are equivalent:
\begin{enumerate}
  \item[{\rm a)}] $\bb{A}$ satisfies (ML).
  \item[{\rm b)}] For every $p\in\bb{Z}$, the projective system
  $(\mathrm{gr}^p(A_k))_{k\in\bb{Z}}$ is essentially constant.
\end{enumerate}
Moreover, when these conditions hold, there is, for every
$p\in\bb{Z}$, a canonical isomorphism
\begin{equation}
\label{0.13.4.3.1}
\tag{13.4.3.1}
  \mathrm{gr}^p(\bb{A})
  \isoto
  \varprojlim_k\mathrm{gr}^p(A_k).
\end{equation}
\end{proposition}

\begin{proof}
It follows at once from \sref{0.13.4.1.2} that \emph{a)} implies
\emph{b)}. The same formula, applied to the projective system
$\bb{A}'$ in the notation of \sref{0.13.4.2}, gives
\sref{0.13.4.3.1} by definition.

For $k\leq h$, put
\[
  A_{kh}=\Im(A_h\longrightarrow A_k).
\]
If $k\leq h\leq j$, then
$A_{kj}\subset A_{kh}\subset A_k$.
Equip $A_{kh}$ with the filtration induced by
$(\mathrm{F}^p(A_k))$.
The transitivity of the morphisms defining $\bb{A}$ shows
immediately that this is also the quotient filtration of
$(\mathrm{F}^p(A_h))$; consequently,
\begin{equation}
\label{0.13.4.3.2}
\tag{13.4.3.2}
  \mathrm{gr}^p(A_{kh})
  =
  \Im\bigl(
    \mathrm{gr}^p(A_h)\longrightarrow\mathrm{gr}^p(A_k)
  \bigr).
\end{equation}

\oldpage[0\textsubscript{III}]{71}
Suppose now that \emph{b)} holds.
For every $p\in\bb{Z}$ and every $k\geq p$, there is an integer
$\mathrm{L}(p,k)$ such that the right-hand side of
\sref{0.13.4.3.2} is constant for $h\geq\mathrm{L}(p,k)$.
Since $\mathrm{gr}^p(A_k)=0$ for $p<k_0$, for fixed $k$ only
finitely many of the integers $\mathrm{L}(p,k)$ need be nonzero as
$p$ ranges over the integers $\leq k$.
Let $\mathrm{L}(k)=m$ be the largest of these integers.
For every $h\geq m$, one has
$A_{kh}\subset A_{km}$, and, by the definition of $m$, the canonical
injection
$A_{kh}\longrightarrow A_{km}$ induces an isomorphism
\[
  \mathrm{gr}^\bullet(A_{kh})
  \simto
  \mathrm{gr}^\bullet(A_{km}).
\]
Since the filtrations are finite, this injection is itself bijective
(Bourbaki, \emph{Alg.\ comm.}, chap.~III, \textsection2,
n\textsuperscript{o}~8, th.~1).
This proves that $\bb{A}$ satisfies (ML).
\end{proof}

\begin{env}[13.4.4]
\label{0.13.4.4}
Suppose that the projective limit
$A=\varprojlim_k A_k$ exists in $\cat{C}$.
In the definitions of \sref{0.13.4.1}, one may then replace $A_k$ by
$A$, and the filtration thereby defined on $A$ still satisfies
\begin{equation}
\label{0.13.4.4.1}
\tag{13.4.4.1}
  \mathrm{gr}^p(A)
  =
  \mathrm{gr}^p(A_p)
  \cap
  \Im(A\longrightarrow A_p).
\end{equation}
\end{env}

\begin{corollary}[13.4.5]
\label{0.13.4.5}
Suppose that $\cat{C}$ is the category of abelian groups.
If the projective system $\bb{A}$ satisfies (ML) and
$A=\varprojlim_k A_k$, then for every $p\in\bb{Z}$ there is a
canonical isomorphism
\begin{equation}
\label{0.13.4.5.1}
\tag{13.4.5.1}
  \mathrm{gr}^p(A)
  \isoto
  \varprojlim_k\mathrm{gr}^p(A_k).
\end{equation}
\end{corollary}

\begin{proof}
Indeed,
\[
  \Im(A_k\longrightarrow A_p)
  =
  \Im(A\longrightarrow A_p)
\]
for all sufficiently large $k$
(Bourbaki, \emph{Top.\ g\'{e}n.}, chap.~II,
3\textsuperscript{e}~\'{e}d., \textsection3,
n\textsuperscript{o}~5, th.~1), and the conclusion follows from
\sref{0.13.4.1.3} and \sref{0.13.4.4.1}.
\end{proof}

\subsection{Projective limits of spectral sequences of filtered complexes}
\label{subsection:0.13.5}

\begin{env}[13.5.1]
\label{0.13.5.1}
Let $\cat{C}$ be an abelian category and let $X^\bullet$ be a complex
of objects of $\cat{C}$ equipped with a filtration
$(\mathrm{F}^p(X^\bullet))_{p\in\bb{Z}}$ such that
$\mathrm{F}^{p_0}(X^\bullet)=X^\bullet$ for some index $p_0$.
For every $k\in\bb{Z}$, consider the complex
\[
  X^\bullet_k=X^\bullet/\mathrm{F}^{k+1}(X^\bullet);
\]
it is canonically equipped with the filtration consisting of
\[
  \mathrm{F}^p(X^\bullet_k)
  =
  \mathrm{F}^p(X^\bullet)/\mathrm{F}^{k+1}(X^\bullet)
  \quad\text{for }p\leq k,
\]
and $\mathrm{F}^p(X^\bullet_k)=0$ for $p\geq k+1$.
Moreover, there are canonical morphisms
$X^\bullet_{k+1}\longrightarrow X^\bullet_k$, which make
$\bb{X}^\bullet=(X^\bullet_k)_{k\in\bb{Z}}$ a \emph{projective system
of filtered complexes of objects of $\cat{C}$}.
We note that this projective system is \emph{strict} and that
$X^\bullet_k=0$ for $k<p_0$.
\end{env}

\begin{env}[13.5.2]
\label{0.13.5.2}
More generally, consider a bounded-below strict projective system
$\bb{X}^\bullet=(X^\bullet_k)_{k\in\bb{Z}}$ of complexes of objects of
$\cat{C}$, and equip each $X^\bullet_k$ with the filtration defined in
\sref{0.13.4.1} (applied in the abelian category of complexes of
objects of $\cat{C}$ with bounded-below filtrations).
The morphisms
$X^\bullet_k\longrightarrow X^\bullet_p$ ($p\leq k$) then become
morphisms of \emph{filtered} complexes with finite filtrations.
The functoriality of the spectral sequences of filtered complexes
(\sref{0.11.2.3}) shows that the transition morphisms of
$\bb{X}^\bullet$ induce morphisms making
\[
  \mathrm{E}(\bb{X}^\bullet)
  =
  \bigl(\mathrm{E}(X^\bullet_k)\bigr)_{k\in\bb{Z}}
\]
a \emph{projective system of spectral sequences}.
\end{env}

\begin{lemma}[13.5.3]
\label{0.13.5.3}
Suppose that the projective system
$\bb{X}^\bullet=(X^\bullet_k)_{k\in\bb{Z}}$ of filtered complexes is
obtained as in \sref{0.13.5.2}. Then:
\begin{enumerate}
  \item[{\rm a)}] For $r\geq p-p_0$, one has
  \[
    \mathrm{B}_r^{pq}(X^\bullet_k)
    =
    \mathrm{B}_\infty^{pq}(X^\bullet_k)
  \]
  for every $k\in\bb{Z}$.
  \item[{\rm b)}] For $k+1\leq p+r$, one has
  \[
    \mathrm{Z}_r^{pq}(X^\bullet_k)
    =
    \mathrm{Z}_\infty^{pq}(X^\bullet_k).
  \]
  \item[{\rm c)}] For $k+1\geq p+r$, the morphisms
  \[
    \mathrm{Z}_r^{pq}(X^\bullet_h)
    \longrightarrow
    \mathrm{Z}_r^{pq}(X^\bullet_k)
    \quad\text{and}\quad
    \mathrm{B}_r^{pq}(X^\bullet_h)
    \longrightarrow
    \mathrm{B}_r^{pq}(X^\bullet_k)
  \]
  are isomorphisms for every $h\geq k$.
\end{enumerate}
\end{lemma}

\oldpage[0\textsubscript{III}]{72}
These three properties follow immediately from the definitions in
\sref{0.11.2.2}, taking into account that
$\mathrm{F}^{p-r+1}(X^\bullet_k)=X^\bullet_k$ when $p-r<p_0$.

\begin{env}[13.5.4]
\label{0.13.5.4}
Suppose that the hypotheses of \sref{0.13.5.3} hold.
Then, for fixed $p$, $q$, and $r$ (with $r$ finite), the projective
systems
\[
  \bigl(\mathrm{Z}_r^{pq}(X^\bullet_k)\bigr)_{k\in\bb{Z}},
  \qquad
  \bigl(\mathrm{B}_r^{pq}(X^\bullet_k)\bigr)_{k\in\bb{Z}},
  \qquad
  \bigl(\mathrm{E}_r^{pq}(X^\bullet_k)\bigr)_{k\in\bb{Z}}
\]
are \emph{essentially constant}; their respective projective limits
will be denoted by
$\mathrm{Z}_r^{pq}(\bb{X}^\bullet)$,
$\mathrm{B}_r^{pq}(\bb{X}^\bullet)$, and
\[
  \mathrm{E}_r^{pq}(\bb{X}^\bullet)
  =
  \mathrm{Z}_r^{pq}(\bb{X}^\bullet)/
  \mathrm{B}_r^{pq}(\bb{X}^\bullet).
\]
The objects $\mathrm{Z}_r^{pq}(\bb{X}^\bullet)$ and
$\mathrm{B}_r^{pq}(\bb{X}^\bullet)$ identify canonically with
subobjects of $\mathrm{E}_1^{pq}(\bb{X}^\bullet)$.
The definition of the $d_r^{pq}$ (M, XV, 1) shows that these
morphisms, taken relative to the $X^\bullet_k$, are also essentially
constant, and hence define morphisms
\begin{equation}
\label{0.13.5.4.1}
\tag{13.5.4.1}
  d_r^{pq}:
  \mathrm{E}_r^{pq}(\bb{X}^\bullet)
  \longrightarrow
  \mathrm{E}_r^{p+r,q-r+1}(\bb{X}^\bullet)
\end{equation}
such that
$d_r^{p+r,q-r+1}\circ d_r^{pq}=0$.
Moreover, there are canonical isomorphisms from
$\Ker(d_r^{pq})$ onto
\[
  \mathrm{Z}_{r+1}^{pq}(\bb{X}^\bullet)/
  \mathrm{B}_r^{pq}(\bb{X}^\bullet)
\]
and from $\Im(d_r^{pq})$ onto
\[
  \mathrm{B}_{r+1}^{p+r,q-r+1}(\bb{X}^\bullet)/
  \mathrm{B}_r^{p+r,q-r+1}(\bb{X}^\bullet).
\]
\end{env}

\begin{lemma}[13.5.5]
\label{0.13.5.5}
Under the hypotheses of \sref{0.13.5.3}, for
$s\geq r>p-p_0$ there is a canonical monomorphism
\begin{equation}
\label{0.13.5.5.1}
\tag{13.5.5.1}
  i:
  \mathrm{E}_s^{pq}(\bb{X}^\bullet)
  \longrightarrow
  \mathrm{E}_r^{pq}(\bb{X}^\bullet)
\end{equation}
and a canonical isomorphism
\begin{equation}
\label{0.13.5.5.2}
\tag{13.5.5.2}
  j_r:
  \mathrm{E}_r^{pq}(\bb{X}^\bullet)
  \simto
  \mathrm{E}_\infty^{pq}(X^\bullet_{p+r-1})
\end{equation}
such that the diagram
\begin{equation}
\label{0.13.5.5.3}
\tag{13.5.5.3}
  \xymatrix{
    \mathrm{E}_s^{pq}(\bb{X}^\bullet)
      \ar[r]^{j_s}\ar[d]_i
    &
    \mathrm{E}_\infty^{pq}(X^\bullet_{p+s-1})
      \ar[d]
    \\
    \mathrm{E}_r^{pq}(\bb{X}^\bullet)
      \ar[r]^{j_r}
    &
    \mathrm{E}_\infty^{pq}(X^\bullet_{p+r-1})
  }
\end{equation}
is commutative, where the right vertical arrow is induced by the
morphism
$X^\bullet_{p+s-1}\longrightarrow X^\bullet_{p+r-1}$.
\end{lemma}

\begin{proof}
The existence of $i$ follows from
\[
  \mathrm{B}_r^{pq}(X^\bullet_k)
  =
  \mathrm{B}_\infty^{pq}(X^\bullet_k)
  \qquad (r>p-p_0)
\]
(\sref{0.13.5.3},~a)).
For $k+1\leq p+r$, one has
\[
  \mathrm{Z}_r^{pq}(X^\bullet_k)
  =
  \mathrm{Z}_\infty^{pq}(X^\bullet_k)
\]
(\sref{0.13.5.3},~b)); in particular,
\[
  \mathrm{Z}_r^{pq}(X^\bullet_{p+r-1})
  =
  \mathrm{Z}_\infty^{pq}(X^\bullet_{p+r-1}).
\]
On the other hand,
$\mathrm{Z}_r^{pq}(X^\bullet_{p+r-1})$ and
$\mathrm{B}_r^{pq}(X^\bullet_{p+r-1})$ identify canonically with
$\mathrm{Z}_r^{pq}(\bb{X}^\bullet)$ and
$\mathrm{B}_r^{pq}(\bb{X}^\bullet)$ by
\sref{0.13.5.3},~c).
This proves the existence of $j_r$ and the commutativity of
\sref{0.13.5.5.3}.
\end{proof}

\begin{corollary}[13.5.6]
\label{0.13.5.6}
Under the hypotheses of \sref{0.13.5.3}, if one of the projective
limits
\[
  \varprojlim_r\mathrm{E}_r^{pq}(\bb{X}^\bullet),
  \qquad
  \varprojlim_k\mathrm{E}_\infty^{pq}(X^\bullet_k)
\]
exists, then so does the other, and there is a canonical isomorphism
\begin{equation}
\label{0.13.5.6.1}
\tag{13.5.6.1}
  j_\infty:
  \varprojlim_r\mathrm{E}_r^{pq}(\bb{X}^\bullet)
  \simto
  \varprojlim_k\mathrm{E}_\infty^{pq}(X^\bullet_k).
\end{equation}
Moreover, the projective system
$(\mathrm{E}_r^{pq}(\bb{X}^\bullet))_{r\in\bb{Z}}$ is essentially
constant (\sref{0.13.4.2}) if and only if the projective system
$(\mathrm{E}_\infty^{pq}(X^\bullet_k))_{k\in\bb{Z}}$ is essentially
constant.
\end{corollary}

\begin{env}[13.5.7]
\label{0.13.5.7}
We denote by
$\mathrm{B}_\infty^{pq}(\bb{X}^\bullet)$ and
$\mathrm{Z}_\infty^{pq}(\bb{X}^\bullet)$ the subobjects of
$\mathrm{E}_1^{pq}(\bb{X}^\bullet)$ equal respectively to
$\mathrm{B}_r^{pq}(\bb{X}^\bullet)$ for $r>p-p_0$ and to
$\inf_r\mathrm{Z}_r^{pq}(\bb{X}^\bullet)$, when the latter exists.
Thus
\[
  \varprojlim_r\mathrm{E}_r^{pq}(\bb{X}^\bullet)
\]
\oldpage[0\textsubscript{III}]{73}
identifies canonically with
\[
  \mathrm{E}_\infty^{pq}(\bb{X}^\bullet)
  =
  \mathrm{Z}_\infty^{pq}(\bb{X}^\bullet)/
  \mathrm{B}_\infty^{pq}(\bb{X}^\bullet).
\]
We note that the objects
$\mathrm{Z}_r^{pq}(\bb{X}^\bullet)$,
$\mathrm{B}_r^{pq}(\bb{X}^\bullet)$,
$\mathrm{E}_r^{pq}(\bb{X}^\bullet)$
($1\leq r\leq+\infty$), and the morphisms $d_r^{pq}$ depend
\emph{functorially} on the projective system $\bb{X}^\bullet$, subject
to the restrictions in \sref{0.13.5.5}, and the morphisms defined in
\sref{0.13.5.5} and \sref{0.13.5.6} are functorial.
\end{env}

\subsection{Spectral sequence of a functor relative to an object equipped with a finite filtration}
\label{subsection:0.13.6}

\begin{env}[13.6.1]
\label{0.13.6.1}
Let $\cat{C}$ and $\cat{C}'$ be two abelian categories, and let
$\mathrm{T}:\cat{C}\longrightarrow\cat{C}'$ be an additive covariant
functor.
Suppose that every object of $\cat{C}$ is isomorphic to a subobject of
an injective object, so that the right derived functors
$\RR^p\mathrm{T}$ ($p\geq0$) exist.
\end{env}

\begin{lemma}[13.6.2]
\label{0.13.6.2}
Let $A$ be an object of $\cat{C}$ equipped with a finite filtration
$(\mathrm{F}^i(A))_{i\in\bb{Z}}$.
There exists an injective resolution
$X^\bullet=(X^j)_{j\geq0}$ of $A$, equipped with a finite filtration
$(\mathrm{F}^i(X^\bullet))_{i\in\bb{Z}}$, such that
\[
  \mathrm{F}^i(A)=A
  \quad\text{(respectively, }\mathrm{F}^i(A)=0\text{)}
\]
implies
\[
  \mathrm{F}^i(X^\bullet)=X^\bullet
  \quad\text{(respectively, }\mathrm{F}^i(X^\bullet)=0\text{)},
\]
and such that, for every $i\in\bb{Z}$,
$\mathrm{F}^i(X^\bullet)$ is an injective resolution of
$\mathrm{F}^i(A)$.
\end{lemma}

\begin{proof}
Let $p$ (respectively, $q>p$) be the largest index such that
$\mathrm{F}^p(A)=A$ (respectively, the smallest index such that
$\mathrm{F}^q(A)=0$).
We argue by induction on $q-p$, the lemma being immediate when
$q-p=1$.
Suppose that an injective resolution ${X''}^\bullet$ of
$A/\mathrm{F}^{q-1}(A)$ having the required properties has been
constructed.
Consider the exact sequence
\[
  0\longrightarrow\mathrm{F}^{q-1}(A)
  \longrightarrow A
  \longrightarrow A/\mathrm{F}^{q-1}(A)
  \longrightarrow0,
\]
take an injective resolution ${X'}^\bullet$ of
$\mathrm{F}^{q-1}(A)$, and then choose an injective resolution
$X^\bullet$ of $A$ so as to obtain an exact sequence
\[
  0\longrightarrow {X'}^\bullet
  \longrightarrow X^\bullet
  \longrightarrow {X''}^\bullet
  \longrightarrow0
\]
compatible with the preceding sequence (M, V, 2.2).
It is clear that $X^\bullet$ has the required properties.
\end{proof}

\begin{corollary}[13.6.3]
\label{0.13.6.3}
Let $B$ be a second object of $\cat{C}$, equipped with a finite
filtration $(\mathrm{F}^i(B))_{i\in\bb{Z}}$, let $s$ be an integer,
and let $u:A\longrightarrow B$ be a morphism such that
\[
  u(\mathrm{F}^i(A))\subset\mathrm{F}^{i+s}(B)
\]
for every $i\in\bb{Z}$.
If $Y^\bullet=(Y^j)_{j\geq0}$ is an injective resolution of $B$
equipped with a filtration
$(\mathrm{F}^i(Y^\bullet))_{i\in\bb{Z}}$ having the properties stated
in \sref{0.13.6.2}, then there exists a morphism
$v:X^\bullet\longrightarrow Y^\bullet$, compatible with $u$, such
that
\[
  v(\mathrm{F}^i(X^\bullet))
  \subset
  \mathrm{F}^{i+s}(Y^\bullet)
\]
for every $i\in\bb{Z}$.
Moreover, any two such morphisms $v$ and $v'$ are homotopic.
\end{corollary}

This follows immediately, by induction on $q-p$, from the preceding
construction and from (M, V, 2.3).

\begin{env}[13.6.4]
\label{0.13.6.4}
Under the hypotheses of \sref{0.13.6.2}, consider the complex
$\mathrm{T}(X^\bullet)$ in $\cat{C}'$.
It is filtered by the complexes
$\mathrm{T}(\mathrm{F}^i(X^\bullet))$, since
$\mathrm{F}^i(X^\bullet)$ is a direct summand of $X^\bullet$.
It follows from \sref{0.13.6.3} that the \emph{spectral sequence} of
this filtered complex depends, up to isomorphism, only on the
filtered object $A$.
Its abutment is the cohomology $\RR^\bullet\mathrm{T}(A)$, equipped
with the filtration
\begin{equation}
\label{0.13.6.4.1}
\tag{13.6.4.1}
  \begin{split}
  \mathrm{F}^p(\RR^n\mathrm{T}(A))
  &=
  \Im\bigl(
    \RR^n\mathrm{T}(\mathrm{F}^p(A))
    \longrightarrow
    \RR^n\mathrm{T}(A)
  \bigr)
  \\
  &=
  \Ker\bigl(
    \RR^n\mathrm{T}(A)
    \longrightarrow
    \RR^n\mathrm{T}(A/\mathrm{F}^p(A))
  \bigr),
  \end{split}
\end{equation}
by \sref{0.11.2.2}, and its $\mathrm{E}_1$-term is
\begin{equation}
\label{0.13.6.4.2}
\tag{13.6.4.2}
  \mathrm{E}_1^{pq}
  =
  \RR^{p+q}\mathrm{T}(\mathrm{gr}^p(A)),
\end{equation}
where, as usual,
$\mathrm{gr}^p(A)=\mathrm{F}^p(A)/\mathrm{F}^{p+1}(A)$.
It is clear from \sref{0.11.2.2} that the filtration of the abutment
is finite and that, for fixed $p$ and $q$, the sequences
\oldpage[0\textsubscript{III}]{74}
\[
  \mathrm{B}_r^{pq}(A)
  =
  \mathrm{B}_r^{pq}(\mathrm{T}(X^\bullet))
  \quad\text{and}\quad
  \mathrm{Z}_r^{pq}(A)
  =
  \mathrm{Z}_r^{pq}(\mathrm{T}(X^\bullet))
\]
are stationary.
The preceding spectral sequence is therefore \emph{biregular}
(\sref{0.11.1.3}).
We denote it by
$\mathrm{E}(A)=(\mathrm{E}_r^{pq}(A))$ and call it the
\emph{spectral sequence of the functor $\mathrm{T}$ relative to the
filtered object $A$}.
\end{env}

\begin{env}[13.6.5]
\label{0.13.6.5}
Now suppose that the hypotheses of \sref{0.13.6.3} hold, and retain
its notation.
Since $\mathrm{F}^i(X^\bullet)$ (respectively,
$\mathrm{F}^i(Y^\bullet)$) is a direct summand of $X^\bullet$
(respectively, $Y^\bullet$), one has
\[
  (\mathrm{T}v)\bigl(\mathrm{T}(\mathrm{F}^i(X^\bullet))\bigr)
  \subset
  \mathrm{T}(\mathrm{F}^{i+s}(Y^\bullet))
\]
for every $i\in\bb{Z}$.
The definitions in \sref{0.11.2.2} then show that, for
$1\leq r\leq+\infty$, $\mathrm{T}v$ defines morphisms
\[
  \mathrm{B}_r^{pq}(\mathrm{T}(X^\bullet))
  \longrightarrow
  \mathrm{B}_r^{p+s,q-s}(\mathrm{T}(Y^\bullet))
\]
and
\[
  \mathrm{Z}_r^{pq}(\mathrm{T}(X^\bullet))
  \longrightarrow
  \mathrm{Z}_r^{p+s,q-s}(\mathrm{T}(Y^\bullet)),
\]
and hence a morphism
\[
  w_r:
  \mathrm{E}_r^{pq}(A)
  \longrightarrow
  \mathrm{E}_r^{p+s,q-s}(B).
\]
Similarly, on the abutments one has morphisms
\[
  u_n:
  \RR^n\mathrm{T}(A)
  \longrightarrow
  \RR^n\mathrm{T}(B)
\]
such that
\[
  u_n\bigl(\mathrm{F}^p(\RR^n\mathrm{T}(A))\bigr)
  \subset
  \mathrm{F}^{p+s}(\RR^n\mathrm{T}(B)).
\]
The definition of the $d_r^{pq}$ (M, XV, 1) also shows that the
diagrams
\[
  \xymatrix{
    \mathrm{E}_r^{pq}(A)
      \ar[r]^{d_r^{pq}}\ar[d]_{w_r}
    &
    \mathrm{E}_r^{p+r,q-r+1}(A)
      \ar[d]^{w_r}
    \\
    \mathrm{E}_r^{p+s,q-s}(B)
      \ar[r]_{d_r^{p+s,q-s}}
    &
    \mathrm{E}_r^{p+r+s,q-r-s+1}(B)
  }
\]
are commutative.
It follows that there is an analogous commutative diagram for the
isomorphisms $\alpha_r^{pq}$, which we leave to the reader to make
explicit.
Finally (\emph{loc.\ cit.}), one also has commutative diagrams on the
abutments:
\[
  \xymatrix{
    \mathrm{E}_\infty^{pq}(A)
      \ar[r]^{\beta^{pq}}\ar[d]_{w_\infty}
    &
    \mathrm{gr}^p(\RR^{p+q}\mathrm{T}(A))
      \ar[d]^{u_{p+q}}
    \\
    \mathrm{E}_\infty^{p+s,q-s}(B)
      \ar[r]_{\beta^{p+s,q-s}}
    &
    \mathrm{gr}^{p+s}(\RR^{p+q}\mathrm{T}(B)).
  }
\]
\end{env}

\begin{env}[13.6.6]
\label{0.13.6.6}
Suppose in particular that there is a ring $S$, equipped with a
\emph{filtration} $(\mathrm{F}^i(S))_{i\in\bb{Z}}$, and a ring
homomorphism
\begin{equation}
\label{0.13.6.6.1}
\tag{13.6.6.1}
  h:S\longrightarrow\Hom_{\cat{C}}(A,A)
\end{equation}
\oldpage[0\textsubscript{III}]{75}
such that, for every $t\in\mathrm{F}^i(S)$,
\[
  h_t(\mathrm{F}^j(A))
  \subset
  \mathrm{F}^{i+j}(A)
\]
for every pair $i,j$.
For brevity, we then say that $A$ is equipped with the structure of a
\emph{filtered $S$--$\cat{C}$-module} over the filtered ring $S$.

Passing to the associated graded objects, every $h_t$, for
$t\in\mathrm{F}^i(S)$, defines a graded endomorphism
$\overline{h}_t$ of $\mathrm{gr}^\bullet(A)$, homogeneous of degree
$i$.
Moreover, this endomorphism depends only on the class of $t$ in
$\mathrm{gr}^i(S)$, and one thereby obtains a homomorphism of
\emph{graded rings}
\[
  \overline{h}:
  \mathrm{gr}^\bullet(S)
  \longrightarrow
  \Hom_{\cat{C}}^\bullet
  \bigl(\mathrm{gr}^\bullet(A),\mathrm{gr}^\bullet(A)\bigr),
\]
where the right-hand side is the ring of graded endomorphisms of
$\mathrm{gr}^\bullet(A)$.
We say that $\mathrm{gr}^\bullet(A)$ is equipped with the structure
of a \emph{graded $\mathrm{gr}^\bullet(S)$--$\cat{C}$-module}.

It follows from \sref{0.13.6.5} that, for
$1\leq r\leq+\infty$, every
$\overline{t}\in\mathrm{gr}^i(S)$ canonically defines, on the
bigraded objects
\[
  \bigl(\mathrm{B}_r^{pq}(A)\bigr)_{(p,q)\in\bb{Z}\times\bb{Z}},
  \qquad
  \bigl(\mathrm{Z}_r^{pq}(A)\bigr)_{(p,q)\in\bb{Z}\times\bb{Z}},
\]
and
\[
  \mathrm{E}_r(A)
  =
  \bigl(\mathrm{E}_r^{pq}(A)\bigr)_{(p,q)\in\bb{Z}\times\bb{Z}},
\]
bigraded endomorphisms of degree $(i,-i)$.
On $\mathrm{E}_r(A)$, for finite $r$, this endomorphism commutes with
the bigraded endomorphism defined by the $d_r^{pq}$.
These endomorphisms satisfy the usual associativity and
distributivity conditions with respect to addition in
$\mathrm{gr}^\bullet(S)$ and in the bigraded objects in question.
For brevity, we say that the latter are
\emph{bigraded $\mathrm{gr}^\bullet(S)$--$\cat{C}'$-modules};
it is immediate that the $\alpha_r^{pq}$ define an isomorphism for
this type of structure.

For every integer $n$, let $\mathrm{B}_r^{(n)}(A)$
(respectively, $\mathrm{Z}_r^{(n)}(A)$ and
$\mathrm{E}_r^{(n)}(A)$) denote the graded subobject of
$\mathrm{B}_r^{\bullet,\bullet}(A)$
(respectively, $\mathrm{Z}_r^{\bullet,\bullet}(A)$ and
$\mathrm{E}_r^{\bullet,\bullet}(A)$) consisting of the terms for
which $p+q=n$ ($1\leq r\leq+\infty$).
It is immediate that these are
\emph{graded $\mathrm{gr}^\bullet(S)$--$\cat{C}'$-modules}.

Finally, every $\overline{t}\in\mathrm{gr}^i(S)$ defines, for every
$n$, a graded endomorphism of degree $i$ on the graded object
$\mathrm{gr}^\bullet(\RR^n\mathrm{T}(A))$, which thereby acquires the
structure of a graded
$\mathrm{gr}^\bullet(S)$--$\cat{C}'$-module.
For $p+q=n$, the $\beta^{pq}$ define an isomorphism
\[
  \mathrm{E}^{(n)}(A)
  \isoto
  \mathrm{gr}^\bullet(\RR^n\mathrm{T}(A))
\]
for this type of structure.

We note that, when $\cat{C}'$ is the category of abelian groups, the
structures of an $S$--$\cat{C}'$-module (respectively, of a graded or
bigraded $\mathrm{gr}^\bullet(S)$--$\cat{C}'$-module) are simply the
usual structures of an $S$-module (respectively, a graded or
bigraded $\mathrm{gr}^\bullet(S)$-module).
\end{env}

\subsection{Derived functors of a projective limit of arguments}
\label{subsection:0.13.7}
\label{0.13.7}

\begin{env}[13.7.1]
\label{0.13.7.1}
Let $\cat{C}$ and $\cat{C}'$ be two abelian categories, and suppose
that every object of $\cat{C}$ is isomorphic to a subobject of an
injective object.
Let
\[
  \mathrm{T}:\cat{C}\longrightarrow\cat{C}'
\]
be an additive covariant functor.
Consider a strict projective system
$\bb{A}=(A_k)_{k\in\bb{Z}}$ in $\cat{C}$, bounded below; more
precisely, suppose that $A_k=0$ for $k<k_0$.
To this system we canonically associate, on each $A_k$, the
filtration $(\mathrm{F}^p(A_k))_{p\in\bb{Z}}$ given by
\sref{0.13.4.1.1}.
Since the projective system is strict, the canonical morphisms
\begin{equation}
\label{0.13.7.1.1}
\tag{13.7.1.1}
  \mathrm{F}^i(A_h)/\mathrm{F}^j(A_h)
  \longrightarrow
  \mathrm{F}^i(A_k)/\mathrm{F}^j(A_k)
  \qquad(h\geq k)
\end{equation}
are isomorphisms for $i\leq j\leq k+1$.
Recall also that
$\mathrm{F}^{k_0}(A_k)=A_k$ and
$\mathrm{F}^{k+1}(A_k)=0$ for every $k$.
\end{env}

\begin{env}[13.7.2]
\label{0.13.7.2}
For each $k$, construct an injective resolution
$\mathrm{X}_k^\bullet=(\mathrm{X}_k^j)_{j\geq0}$ of $A_k$ having the
properties of \sref{0.13.6.2}.
The canonical morphisms $A_{k+1}\longrightarrow A_k$ allow us, by
\sref{0.13.6.3}, to define for each $k$ a morphism of complexes
\[
  \mathrm{X}_{k+1}^\bullet\longrightarrow\mathrm{X}_k^\bullet
\]
compatible with the filtrations and making
$\bb{X}^\bullet=(\mathrm{X}_k^\bullet)_{k\in\bb{Z}}$ a projective
system of complexes.
\oldpage[0\textsubscript{III}]{76}
We may moreover suppose that this projective system is strict.
Indeed, by \sref{0.13.7.1.1},
\[
  A_k
  \simeq
  A_{k+1}/\mathrm{F}^{k+1}(A_{k+1}).
\]
Thus, in constructing $\mathrm{X}_{k+1}^\bullet$, we may take the
injective resolution of
$A_{k+1}/\mathrm{F}^{k+1}(A_{k+1})$ to be
\emph{equal} to $\mathrm{X}_k^\bullet$.
It then follows from (M, V, 2.3) that the morphism of complexes
$\mathrm{X}_{k+1}^\bullet\longrightarrow\mathrm{X}_k^\bullet$ may be
constructed so that, on passing to quotients, it induces an
isomorphism
\[
  \mathrm{X}_{k+1}^\bullet/
  \mathrm{F}^{k+1}(\mathrm{X}_{k+1}^\bullet)
  \simto
  \mathrm{X}_k^\bullet
\]
respecting the filtrations; this is precisely the condition of
\sref{0.13.5.1}.
\end{env}

\begin{env}[13.7.3]
\label{0.13.7.3}
By construction, the projective system
$\mathrm{T}(\bb{X}^\bullet)$ of complexes
$\mathrm{T}(\mathrm{X}_k^\bullet)$ satisfies the hypotheses of
\sref{0.13.5.3}.
The results of \sref{0.13.5.4}, \sref{0.13.5.5}, and
\sref{0.13.5.6} therefore apply to the spectral sequences
\[
  \mathrm{E}\bigl(\mathrm{T}(\mathrm{X}_k^\bullet)\bigr)
  =
  \mathrm{E}(A_k).
\]
For $1\leq r\leq+\infty$, we write
$\mathrm{E}_r^{pq}(\bb{A})$ in place of
$\mathrm{E}_r^{pq}(\mathrm{T}(\bb{X}^\bullet))$
(cf.\ \sref{0.13.5.7} when $r=+\infty$), and use the same convention
for the analogous notation.
In particular,
\begin{equation}
\label{0.13.7.3.1}
\tag{13.7.3.1}
  \mathrm{E}_1^{pq}(\bb{A})
  =
  \RR^{p+q}\mathrm{T}\bigl(\mathrm{gr}^p(\bb{A})\bigr),
\end{equation}
by \sref{0.13.6.4.2} and the fact that the projective system
$(\mathrm{gr}^p(A_k))_{k\in\bb{Z}}$ is essentially constant.
These results, together with \sref{0.13.4.3}, give the following
proposition, first proved by Shih Weishu by a different method
(unpublished).
\end{env}

\begin{proposition}[13.7.4]
\label{0.13.7.4}
\emph{(Shih).}
Let $n$ be an integer.
The following two conditions are equivalent:
\begin{enumerate}
  \item[{\rm a)}] For every pair $(p,q)$ such that $p+q=n$, the
  projective system
  $(\mathrm{E}_r^{pq}(\bb{A}))_{r\geq2}$ is essentially constant.
  \item[{\rm b)}] The projective system
  \[
    \RR^n\mathrm{T}(\bb{A})
    =
    \bigl(\RR^n\mathrm{T}(A_k)\bigr)_{k\in\bb{Z}}
  \]
  satisfies {\rm(ML)}.
\end{enumerate}
Moreover, when these conditions hold, there is a canonical
isomorphism
\begin{equation}
\label{0.13.7.4.1}
\tag{13.7.4.1}
  \mathrm{gr}^p\bigl(\RR^n\mathrm{T}(\bb{A})\bigr)
  \isoto
  \mathrm{E}_\infty^{p,n-p}(\bb{A})
  \qquad(p\in\bb{Z}).
\end{equation}
\end{proposition}

\begin{proof}
By \sref{0.13.5.6}, condition {\rm a)} says precisely that, for
$p+q=n$, the projective system
$(\mathrm{E}_\infty^{pq}(A_k))_{k\in\bb{Z}}$ is essentially
constant.
On the other hand,
$\mathrm{gr}^p(\RR^n\mathrm{T}(A_k))$ is canonically isomorphic to
$\mathrm{E}_\infty^{p,n-p}(A_k)$.
It follows from \sref{0.13.4.3} that {\rm a)} and {\rm b)} are
equivalent, and \sref{0.13.7.4.1} is just
\sref{0.13.5.6.1} in the case under consideration.
\end{proof}

\begin{corollary}[13.7.5]
\label{0.13.7.5}
Let $\sh{F}=(\sh{F}_k)_{k\in\bb{N}}$ be a projective system of
sheaves of abelian groups satisfying conditions {\rm(i)}, {\rm(ii)},
and {\rm(iii)} of \sref{0.13.3.1}, and put
$\sh{F}=\varprojlim_k\sh{F}_k$.
Suppose that, for the functor
$\sh{G}\longmapsto\Gamma(X,\sh{G})$, the projective system
$(\mathrm{E}_r^{pq}(\sh{F}))_r$ is essentially constant whenever
$p+q=n$ or $p+q=n+1$.
On $\HH^{n+1}(X,\sh{F})$, consider the filtration
\[
  \mathrm{F}^p\bigl(\HH^{n+1}(X,\sh{F})\bigr)
  =
  \Ker\bigl(
    \HH^{n+1}(X,\sh{F})
    \longrightarrow
    \HH^{n+1}(X,\sh{F}_{p-1})
  \bigr).
\]
Then there is a canonical isomorphism
\begin{equation}
\label{0.13.7.5.1}
\tag{13.7.5.1}
  \mathrm{gr}^p\bigl(\HH^{n+1}(X,\sh{F})\bigr)
  \isoto
  \mathrm{E}_\infty^{p,n-p+1}(\sh{F})
  \qquad(p\in\bb{Z}).
\end{equation}
\end{corollary}

\begin{proof}
\oldpage[0\textsubscript{III}]{77}
Apply \sref{0.13.7.4} with $\cat{C}$ the category of sheaves of
abelian groups on $X$, $\cat{C}'$ the category of abelian groups,
and $\mathrm{T}=\Gamma$.
It gives, for every $p\in\bb{Z}$, a canonical isomorphism
\[
  \mathrm{gr}^p\bigl(\RR^{n+1}\Gamma(\sh{F})\bigr)
  \isoto
  \mathrm{E}_\infty^{p,n-p+1}(\sh{F}).
\]
Moreover, by \sref{0.13.7.4}, the projective system
$(\HH^n(X,\sh{F}_k))_{k\in\bb{Z}}$ satisfies {\rm(ML)}; hence
\sref{0.13.3.1} gives a canonical isomorphism
\[
  \HH^{n+1}(X,\sh{F})
  \isoto
  \varprojlim_k\HH^{n+1}(X,\sh{F}_k).
\]

Since the projective system
$\RR^{n+1}\Gamma(\sh{F})$ satisfies {\rm(ML)} by
\sref{0.13.7.4}, there is a canonical isomorphism
\[
  \mathrm{gr}^p\bigl(\RR^{n+1}\Gamma(\sh{F})\bigr)
  \isoto
  \varprojlim_k
  \mathrm{gr}^p\bigl(\HH^{n+1}(X,\sh{F}_k)\bigr)
\]
by \sref{0.13.4.3}, and a canonical isomorphism
\[
  \varprojlim_k
  \mathrm{gr}^p\bigl(\HH^{n+1}(X,\sh{F}_k)\bigr)
  \isoto
  \mathrm{gr}^p\left(
    \varprojlim_k\HH^{n+1}(X,\sh{F}_k)
  \right)
\]
by \sref{0.13.4.5}.
It remains only to check that the preceding isomorphism is
compatible with the filtrations on its two sides.
This follows immediately from the definitions and, for every $p$,
from the commutativity of
\[
  \xymatrix{
    \HH^{n+1}(X,\sh{F})
      \ar[r]^-{\sim}\ar[dr]
    &
    \varprojlim_k\HH^{n+1}(X,\sh{F}_k)
      \ar[d]
    \\
    &
    \HH^{n+1}(X,\sh{F}_{p-1}).
  }
\]
\end{proof}

\begin{env}[13.7.6]
\label{0.13.7.6}
Let $S$ be a ring equipped with a filtration
$(\mathrm{F}^i(S))_{i\in\bb{Z}}$ such that
$\mathrm{F}^0(S)=S$ (and hence
$\mathrm{gr}^i(S)=0$ for $i<0$).
Suppose that each $A_k$, with the filtration defined in
\sref{0.13.7.1}, is a filtered
$S$--$\cat{C}$-module in the sense of \sref{0.13.6.6}, and that the
morphisms $A_h\longrightarrow A_k$ for $k\leq h$ are morphisms of
filtered $S$--$\cat{C}$-modules.
For brevity, we then say that $\bb{A}$ is a
\emph{projective system of filtered $S$--$\cat{C}$-modules}.

It is immediate that, for $k\leq h$ and
$1\leq r\leq+\infty$, the morphisms
\[
  \mathrm{B}_r^{pq}(A_h)\longrightarrow\mathrm{B}_r^{pq}(A_k),
  \qquad
  \mathrm{Z}_r^{pq}(A_h)\longrightarrow\mathrm{Z}_r^{pq}(A_k)
\]
are morphisms of bigraded
$\mathrm{gr}^\bullet(S)$--$\cat{C}'$-modules
(\sref{0.13.6.5}), and that, for finite $r$, the families
\[
  \bigl(\mathrm{Z}_r^{pq}(\bb{A})\bigr),\qquad
  \bigl(\mathrm{B}_r^{pq}(\bb{A})\bigr),\qquad
  \bigl(\mathrm{E}_r^{pq}(\bb{A})\bigr)
\]
are bigraded
$\mathrm{gr}^\bullet(S)$--$\cat{C}'$-modules, the first two being
submodules of the third.
We again write
$\mathrm{Z}_r^{(n)}(\bb{A})$,
$\mathrm{B}_r^{(n)}(\bb{A})$, and
$\mathrm{E}_r^{(n)}(\bb{A})$ for the respective graded subobjects
obtained by retaining only the terms with $p+q=n$; these are graded
$\mathrm{gr}^\bullet(S)$--$\cat{C}'$-modules.

When the projective system
$(\mathrm{E}_r^{pq}(\bb{A}))_r$ is essentially constant,
$\mathrm{E}_\infty^{pq}(\bb{A})$ is likewise a bigraded
$\mathrm{gr}^\bullet(S)$--$\cat{C}'$-module, and each
$\mathrm{E}_\infty^{(n)}(\bb{A})$ is a graded
$\mathrm{gr}^\bullet(S)$--$\cat{C}'$-module.
Moreover, for each $k$, the isomorphisms
\[
  \beta^{p,n-p}:
  \mathrm{E}_\infty^{p,n-p}(A_k)
  \isoto
  \mathrm{gr}^p\bigl(\RR^n\mathrm{T}(A_k)\bigr)
\]
form an isomorphism of graded
$\mathrm{gr}^\bullet(S)$--$\cat{C}'$-modules from
$\mathrm{E}_\infty^{(n)}(A_k)$ onto
$\mathrm{gr}^\bullet(\RR^n\mathrm{T}(A_k))$.
Under the preceding hypotheses, the isomorphism
\[
  \beta^{p,n-p}:
  \mathrm{E}_\infty^{(n)}(\bb{A})
  \isoto
  \varprojlim_k
  \mathrm{gr}^\bullet\bigl(\RR^n\mathrm{T}(A_k)\bigr)
\]
therefore also respects these structures.
The same is plainly true of the canonical isomorphism
\[
  \mathrm{gr}^\bullet\bigl(\RR^n\mathrm{T}(\bb{A})\bigr)
  \isoto
  \varprojlim_k
  \mathrm{gr}^\bullet\bigl(\RR^n\mathrm{T}(A_k)\bigr);
\]
consequently, the isomorphisms \sref{0.13.7.4.1} form an
isomorphism of graded
$\mathrm{gr}^\bullet(S)$--$\cat{C}'$-modules.
\end{env}

\begin{proposition}[13.7.7]
\label{0.13.7.7}
\oldpage[0\textsubscript{III}]{78}
Let $S$ be a Noetherian $\mathfrak{J}$-adic ring.
Suppose that $\cat{C}$ is an abelian category in which every object
is isomorphic to a subobject of an injective object, and let
$\mathrm{T}$ be an additive covariant functor from $\cat{C}$ to the
category of abelian groups.
Let $\bb{A}=(A_k)_{k\in\bb{Z}}$ be a strict projective system of
filtered $S$--$\cat{C}$-modules, for the $\mathfrak{J}$-adic
filtration on $S$, bounded below.
For a fixed integer $n$, suppose that the following condition holds:
\begin{itemize}
  \item[$(\mathrm{F}_n)$]
  \emph{The graded $\mathrm{gr}^\bullet(S)$-module}
  \[
    \mathrm{E}_1^{(m)}(\bb{A})
    =
    \bigl(
      \RR^m\mathrm{T}(\mathrm{gr}^p(\bb{A}))
    \bigr)_{p\in\bb{Z}}
  \]
  \emph{of \sref{0.13.7.3.1} is of finite type for
  $m=n$ and $m=n+1$.}
\end{itemize}
Then:
\begin{enumerate}
  \item[{\rm(i)}] The projective systems
  $(\RR^n\mathrm{T}(A_k))_{k\in\bb{Z}}$ and
  $(\RR^{n+1}\mathrm{T}(A_k))_{k\in\bb{Z}}$ satisfy {\rm(ML)}.
  \item[{\rm(ii)}] If
  \[
    \RR'^{\,n}\mathrm{T}(\bb{A})
    =
    \varprojlim_k\RR^n\mathrm{T}(A_k),
  \]
  then $\RR'^{\,n}\mathrm{T}(\bb{A})$ is an $S$-module of finite
  type.
  \item[{\rm(iii)}] The filtration on
  $\RR'^{\,n}\mathrm{T}(\bb{A})$ defined by
  \[
    \mathrm{F}^p\bigl(\RR'^{\,n}\mathrm{T}(\bb{A})\bigr)
    =
    \Ker\bigl(
      \RR'^{\,n}\mathrm{T}(\bb{A})
      \longrightarrow
      \RR^n\mathrm{T}(A_{p-1})
    \bigr)
    \qquad(p\in\bb{Z})
  \]
  is $\mathfrak{J}$-good; that is,
  \[
    \mathfrak{J}\,
    \mathrm{F}^p\bigl(\RR'^{\,n}\mathrm{T}(\bb{A})\bigr)
    \subset
    \mathrm{F}^{p+1}\bigl(\RR'^{\,n}\mathrm{T}(\bb{A})\bigr)
  \]
  for every $p$, with equality for all sufficiently large $p$.
  In particular, the topology on
  $\RR'^{\,n}\mathrm{T}(\bb{A})$ defined by this filtration is the
  $\mathfrak{J}$-adic topology.
  \item[{\rm(iv)}] The projective system
  $(\mathrm{E}_r^{pq}(\bb{A}))_r$ is essentially constant whenever
  $p+q=n$ or $p+q=n+1$.
  Thus $\mathrm{E}_\infty^{pq}(\bb{A})$ is defined
  (\sref{0.13.5.7}), and the isomorphisms
  \begin{equation}
  \label{0.13.7.7.1}
  \tag{13.7.7.1}
    \mathrm{gr}^p\bigl(\RR'^{\,n}\mathrm{T}(\bb{A})\bigr)
    \isoto
    \mathrm{E}_\infty^{p,n-p}(\bb{A})
    \qquad(p\in\bb{Z})
  \end{equation}
  form an isomorphism of graded
  $\mathrm{gr}^\bullet(S)$-modules.
\end{enumerate}
\end{proposition}

\begin{proof}
In view of \sref{0.13.7.7.1}, and taking account of the
isomorphisms \sref{0.13.7.4.1}, we shall by abuse of notation also
write $\RR^n\mathrm{T}(\bb{A})$ for the projective limit
$\RR'^{\,n}\mathrm{T}(\bb{A})$ of the projective system
$(\RR^n\mathrm{T}(A_k))$.

The graded ring $\mathrm{gr}^\bullet(S)$ is Noetherian
(Bourbaki, \emph{Alg.\ comm.}, chap.~III, \textsection2,
n\textsuperscript{o}~9, cor.~5 of th.~2).
Hence the increasing sequence of graded
$\mathrm{gr}^\bullet(S)$-submodules
$\mathrm{B}_r^{(m)}(\bb{A})$ of
$\mathrm{E}_1^{(m)}(\bb{A})$ is stationary for $m=n$ and
$m=n+1$.
Condition {\rm(b)} of \sref{0.11.1.10} is therefore satisfied.
It follows that condition {\rm(a)} of \sref{0.13.7.4} holds for
$n$ and for $n+1$, proving {\rm(i)}.

Furthermore, the isomorphisms \sref{0.13.7.4.1}, together with the
remarks of \sref{0.13.7.6}, show that
$\mathrm{gr}^\bullet(\RR^n\mathrm{T}(\bb{A}))$ is a graded
$\mathrm{gr}^\bullet(S)$-module isomorphic to
\[
  \mathrm{E}_\infty^{(n)}(\bb{A})
  =
  \mathrm{Z}_\infty^{(n)}(\bb{A})
  /
  \mathrm{B}_\infty^{(n)}(\bb{A}).
\]
Since $\mathrm{Z}_\infty^{(n)}(\bb{A})$ is a submodule of
$\mathrm{E}_1^{(n)}(\bb{A})$, it is of finite type, and so is
$\mathrm{E}_\infty^{(n)}(\bb{A})$.
Finally, for the filtration
$(\mathrm{F}^p(\RR'^{\,n}\mathrm{T}(\bb{A})))$, it follows from
\sref{0.13.4.5} that
\[
  \mathrm{gr}^\bullet\bigl(\RR^n\mathrm{T}(\bb{A})\bigr)
  \quad\hbox{and}\quad
  \mathrm{gr}^\bullet\bigl(\RR'^{\,n}\mathrm{T}(\bb{A})\bigr)
\]
are isomorphic graded $\mathrm{gr}^\bullet(S)$-modules.
This proves {\rm(iv)}.
Assertions {\rm(ii)} and {\rm(iii)} will follow from these results
and the following lemma.
\end{proof}

\begin{lemma}[13.7.7.2]
\label{0.13.7.7.2}
Let $S$ be a Noetherian $\mathfrak{J}$-adic ring, and let $M$
be an $S$-module equipped with a codiscrete filtration
$(\mathrm{F}^p(M))_{p\in\bb{Z}}$ such that
\[
  \mathfrak{J}\mathrm{F}^p(M)
  \subset
  \mathrm{F}^{p+1}(M).
\]
Thus $M$ is a filtered module over $S$ endowed with its
$\mathfrak{J}$-adic filtration.
Suppose moreover that $M$ is separated for the topology defined by
$(\mathrm{F}^p(M))$.
Then the following conditions are equivalent:
\begin{enumerate}
  \item[{\rm a)}] $M$ is an $S$-module of finite type and
  $(\mathrm{F}^p(M))$ is a $\mathfrak{J}$-good filtration.
  \item[{\rm b)}] $\mathrm{gr}^\bullet(M)$ is a
  $\mathrm{gr}^\bullet(S)$-module of finite type.
  \item[{\rm c)}] The $\mathrm{gr}^p(M)$ are $S$-modules of
  finite type and, for all sufficiently large $p$, the canonical
  homomorphisms
  \begin{equation}
  \label{0.13.7.7.3}
  \tag{13.7.7.3}
    \mathfrak{J}\otimes_S\mathrm{gr}^p(M)
    \longrightarrow
    \mathrm{gr}^{p+1}(M)
  \end{equation}
  \oldpage[0\textsubscript{III}]{79}
  are surjective.
  They are induced by
  $\mathfrak{J}\otimes_S\mathrm{F}^p(M)
  \longrightarrow\mathrm{F}^{p+1}(M)$, taking account of the fact
  that the image of the composite homomorphism
  \[
    \mathfrak{J}\otimes_S\mathrm{F}^{p+1}(M)
    \longrightarrow
    \mathfrak{J}\otimes_S\mathrm{F}^p(M)
    \longrightarrow
    \mathrm{F}^{p+1}(M)
  \]
  is
  $\mathfrak{J}\mathrm{F}^{p+1}(M)
  \subset\mathrm{F}^{p+2}(M)$.
\end{enumerate}

For the proof, see Bourbaki, \emph{Alg.\ comm.}, chap.~III,
\textsection3, n\textsuperscript{o}~1, prop.~3.
\end{lemma}

\begin{env}[13.7.7.4]
\label{0.13.7.7.4}
To apply the lemma \sref{0.13.7.7.2}, it remains to observe that the
topology on $\RR'^{\,n}\mathrm{T}(\bb{A})$ defined by the filtration
under consideration makes it a separated and complete $S$-module:
this is the projective-limit topology of the discrete groups
$\RR^n\mathrm{T}(A_k)$.
On the other hand, if $A_k=0$ for $k<k_0$, then also
$\RR^n\mathrm{T}(A_k)=0$ for $k<k_0$; hence
\[
  \mathrm{F}^{k_0}\bigl(\RR'^{\,n}\mathrm{T}(\bb{A})\bigr)
  =
  \RR'^{\,n}\mathrm{T}(\bb{A}).
\]
Thus all the hypotheses of the lemma are indeed satisfied.
\end{env}

\begin{corollary}[13.7.8]
\label{0.13.7.8}
If hypothesis $(\mathrm{F}_n)$ holds, then for every $f\in S$ there
is a canonical isomorphism
\begin{equation}
\label{0.13.7.8.1}
\tag{13.7.8.1}
  \varprojlim_k\bigl(\RR^n\mathrm{T}(A_k)\bigr)_f
  \isoto
  \RR^n\mathrm{T}(\bb{A})\otimes_S S_{(f)}.
\end{equation}
\end{corollary}

\begin{proof}
Indeed, $\RR^n\mathrm{T}(\bb{A})$ is an $S$-module of finite type,
and $S_{(f)}$ is a Noetherian adic $S$-algebra
(\sref[0\textsubscript{I}]{0.7.6.11}), namely the separated
completion of $S_f$ for the $\mathfrak{J}$-preadic topology
(\sref[0\textsubscript{I}]{0.7.6.2}).
It follows from \sref[0\textsubscript{I}]{0.7.7.8} and
\sref[0\textsubscript{I}]{0.7.7.1} that
\[
  \RR^n\mathrm{T}(\bb{A})\otimes_S S_{(f)}
\]
is isomorphic to the separated completion of
$\RR^n\mathrm{T}(\bb{A})\otimes_S S_f$ for the
$\mathfrak{J}$-preadic topology.
A fundamental system of neighbourhoods of zero for this topology
is
\[
  \bigl(
    \mathfrak{J}^p\RR^n\mathrm{T}(\bb{A})
  \bigr)\otimes_S S_f;
\]
the groups
\[
  \mathrm{F}^p\bigl(\RR^n\mathrm{T}(\bb{A})\bigr)
  \otimes_S S_f
\]
therefore also form such a system.
The latter group is the kernel of the canonical map
\[
  \bigl(\RR^n\mathrm{T}(\bb{A})\bigr)_f
  \longrightarrow
  \bigl(\RR^n\mathrm{T}(A_{p-1})\bigr)_f.
\]
Consequently, the separated group associated with
$\RR^n\mathrm{T}(\bb{A})\otimes_S S_f$ identifies with a subgroup
$G$ of
\[
  \varprojlim_k\bigl(\RR^n\mathrm{T}(A_k)\bigr)_f.
\]
But the projective system
$((\RR^n\mathrm{T}(A_k))_f)_k$ plainly satisfies {\rm(ML)}, and the
image of $(\RR^n\mathrm{T}(\bb{A}))_f$ in each
$(\RR^n\mathrm{T}(A_k))_f$ equals the common image of the
$(\RR^n\mathrm{T}(A_h))_f$ for all sufficiently large $h\geq k$.
It follows at once that $G$ is \emph{everywhere dense} in
$\varprojlim_k(\RR^n\mathrm{T}(A_k))_f$.
Since the latter group is separated and complete, the corollary is
proved.
\end{proof}
